\documentclass[10pt]{article}
\usepackage{amsmath,amssymb}
\usepackage{longtable}

\usepackage[a4paper,left=2cm,right=2cm,top=1.5cm,bottom=2cm]{geometry}
\usepackage{xcolor}
\setlength{\parindent}{0pt}
\setlength{\emergencystretch}{2em}
\begin{document}

\title{Exercise: AdaBoost and Gradient Boosting for Heart Disease Diagnosis - Solutions}
\author{}
\date{}
\maketitle
\subsection*{Learning objectives}
By the end of this exercise, you should be able to:

\begin{itemize}
\item
Explain what a \textbf{\textit{decision stump}} is.

\item
Understand how a stump makes a simple classification.

\item
Explain why \textbf{\textit{AdaBoost assigns weights to training examples}}.

\item
Calculate the error and weight of a weak learner in AdaBoost.

\item
Update patient weights after a weak learner makes predictions.

\item
Understand how several weak learners are combined into a strong classifier.

\item
Explain the basic idea behind \textbf{\textit{gradient boosting}}.

\item
Understand how gradient boosting focuses subsequent models on the errors made by previous models.

\item
Distinguish the way \textbf{\textit{AdaBoost}} and \textbf{\textit{gradient boosting}} build ensembles.

\end{itemize}
\section*{Clinical scenario}
A hospital wants to develop a simple machine-learning system to help diagnose \textbf{\textit{heart disease}}.

For each patient, four characteristics are recorded:

\begin{enumerate}
\item
\textbf{\textit{Age}} - patient's age in years.

\item
\textbf{\textit{Body weight}} - patient's weight in kg.

\item
\textbf{\textit{Blocked arteries}} - whether the patient has blocked arteries (\texttt{Yes}/\texttt{No}).

\item
\textbf{\textit{Chest pain}} - whether the patient experiences characteristic chest pain (\texttt{Yes}/\texttt{No}).

\end{enumerate}
A cardiologist has already determined whether each patient actually has heart disease. This diagnosis will be treated as the \textbf{\textit{target variable}}.

For this exercise:

\begin{itemize}
\item
\texttt{+1} = heart disease

\item
\texttt{-1} = no heart disease

\end{itemize}
\textbf{\textit{Important:}} This is a simplified educational dataset. The variables and values are fictional and are not intended to represent a real clinical diagnostic system.

\section*{The 10 patients}
\begingroup
\small
\setlength{\tabcolsep}{4pt}
\begin{longtable}{l|l|l|l|l|l}
Patient & Age & Weight (kg) & \shortstack[l]{Blocked \\ arteries} & Chest pain & \shortstack[l]{Actual \\ diagnosis} \\[1.5pt]
\hline
\endfirsthead
Patient & Age & Weight (kg) & \shortstack[l]{Blocked \\ arteries} & Chest pain & \shortstack[l]{Actual \\ diagnosis} \\[1.5pt]
\hline
\endhead
1 & 25 & 62 & No & No & -1 \\[1.5pt]
2 & 35 & 72 & No & Yes & +1 \\[1.5pt]
3 & 45 & 85 & Yes & No & +1 \\[1.5pt]
4 & 50 & 95 & No & Yes & +1 \\[1.5pt]
5 & 55 & 80 & Yes & Yes & +1 \\[1.5pt]
6 & 60 & 90 & Yes & No & +1 \\[1.5pt]
7 & 65 & 70 & No & No & -1 \\[1.5pt]
8 & 40 & 100 & No & No & -1 \\[1.5pt]
9 & 70 & 88 & Yes & No & +1 \\[1.5pt]
10 & 30 & 110 & No & No & -1 \\[1.5pt]
\end{longtable}
\endgroup
There are \textbf{\textit{6 patients with heart disease}} and \textbf{\textit{4 without heart disease}}.

\clearpage
\section*{Part I - Decision stumps}
A \textbf{\textit{decision stump}} is a decision tree with only \textbf{\textit{one split}}.

For example:

If age > 47, predict heart disease; otherwise predict no heart disease.

This is a very simple model. It uses only \textbf{\textit{one feature and one threshold}}.

For categorical variables, a stump might instead be:

If blocked arteries = Yes, predict heart disease; otherwise predict no heart disease.

\subsection*{Task 1.1 - Build some stumps}
Consider the following possible stumps.

\subsubsection*{Stump A - Age}
\[
\text{If Age}>47:
\quad \hat y=+1
\]

Otherwise:

\[
\hat y=-1
\]

\subsubsection*{Stump B - Blocked arteries}
\[
\text{If Blocked arteries = Yes}:
\quad \hat y=+1
\]

Otherwise:

\[
\hat y=-1
\]

\subsubsection*{Stump C - Chest pain}
\[
\text{If Chest pain = Yes}:
\quad \hat y=+1
\]

Otherwise:

\[
\hat y=-1
\]

\subsubsection*{Exercise 1.1}
For each stump:

\begin{enumerate}
\item
Predict the diagnosis for all 10 patients.

\item
Count the number of incorrect predictions.

\item
Calculate the classification error:

\end{enumerate}
\[
\text{Error}
=
\frac{\text{number of incorrect predictions}}{10}
\]

Complete the table:

\begingroup
\small
\setlength{\tabcolsep}{4pt}
\begin{longtable}{l|l|l|l}
Stump & Feature used & Number incorrect & Error \\[1.5pt]
\hline
\endfirsthead
Stump & Feature used & Number incorrect & Error \\[1.5pt]
\hline
\endhead
\rule{0pt}{2.4ex}A & Age &  $\phantom{-0.00}$ &  $\phantom{-0.00}$ \\[3pt]
\rule{0pt}{2.4ex}B & Blocked arteries &  $\phantom{-0.00}$ &  $\phantom{-0.00}$ \\[3pt]
\rule{0pt}{2.4ex}C & Chest pain &  $\phantom{-0.00}$ &  $\phantom{-0.00}$ \\[3pt]
\end{longtable}
\endgroup
\begingroup\color{red}

\subsubsection*{Solution}
\subsubsection*{Stump A - Age}
Rule:

\[
\text{Age}>47 \Rightarrow +1
\]

Patients predicted positive: \textbf{\textit{4, 5, 6, 7, 9}}

\begingroup
\small
\setlength{\tabcolsep}{4pt}
\begin{longtable}{l|l|l|l}
Patient & Actual & Prediction & Correct? \\[1.5pt]
\hline
\endfirsthead
Patient & Actual & Prediction & Correct? \\[1.5pt]
\hline
\endhead
1 & -1 & -1 & x \\[1.5pt]
2 & +1 & -1 & o \\[1.5pt]
3 & +1 & -1 & o \\[1.5pt]
4 & +1 & +1 & x \\[1.5pt]
5 & +1 & +1 & x \\[1.5pt]
6 & +1 & +1 & x \\[1.5pt]
7 & -1 & +1 & o \\[1.5pt]
8 & -1 & -1 & x \\[1.5pt]
9 & +1 & +1 & x \\[1.5pt]
10 & -1 & -1 & x \\[1.5pt]
\end{longtable}
\endgroup
There are \textbf{\textit{3 errors}}, so

\[
{\epsilon_A=0.30}
\]

\textbf{\textit{Correction:}} The original exercise's implied error for Stump A should therefore be 0.30, not 0.10.

\subsubsection*{Stump B - Blocked arteries}
Rule:

\[
\text{Blocked arteries = Yes} \Rightarrow +1
\]

Patients 3, 5, 6, and 9 are predicted positive.

Patients 2 and 4 are infected but have no blocked arteries, so they are incorrectly classified.

Thus:

\[
{\epsilon_B=\frac{2}{10}=0.20}
\]

\subsubsection*{Stump C - Chest pain}
Rule:

\[
\text{Chest pain = Yes} \Rightarrow +1
\]

Patients 2, 4, and 5 are correctly classified as positive.

Patients 3, 6, and 9 have heart disease but no chest pain, so they are incorrectly classified.

Thus:

\[
{\epsilon_C=\frac{3}{10}=0.30}
\]

\subsubsection*{Summary}
\begingroup
\small
\setlength{\tabcolsep}{4pt}
\begin{longtable}{l|l|l|l}
Stump & Feature & Errors & Error \\[1.5pt]
\hline
\endfirsthead
Stump & Feature & Errors & Error \\[1.5pt]
\hline
\endhead
A & Age & 3 & 0.30 \\[1.5pt]
B & Blocked arteries & 2 & 0.20 \\[1.5pt]
C & Chest pain & 3 & 0.30 \\[1.5pt]
\end{longtable}
\endgroup
Therefore, among these three candidate stumps, \textbf{\textit{Stump B is selected}} because it has the smallest error.

\par\endgroup

\clearpage
\section*{Part II - Which stump should AdaBoost choose?}
AdaBoost starts by giving \textbf{\textit{every training example the same weight}}.

Because there are 10 patients:

\[
w_i=\frac{1}{10}=0.1
\]

for every patient.

Thus:

\begingroup
\small
\setlength{\tabcolsep}{4pt}
\begin{longtable}{l|l}
Patient & Initial weight \\[1.5pt]
\hline
\endfirsthead
Patient & Initial weight \\[1.5pt]
\hline
\endhead
1 & 0.10 \\[1.5pt]
2 & 0.10 \\[1.5pt]
3 & 0.10 \\[1.5pt]
4 & 0.10 \\[1.5pt]
5 & 0.10 \\[1.5pt]
6 & 0.10 \\[1.5pt]
7 & 0.10 \\[1.5pt]
8 & 0.10 \\[1.5pt]
9 & 0.10 \\[1.5pt]
10 & 0.10 \\[1.5pt]
\end{longtable}
\endgroup
AdaBoost calculates the \textbf{\textit{weighted error}} of each stump:

\[
\epsilon
=
\sum_{i=1}^{10}
w_i I(y_i\neq h(x_i))
\]

where $I(\cdot)$ equals 1 when the prediction is wrong and 0 otherwise.

\subsubsection*{Exercise 2.1}
\begin{enumerate}
\item
Calculate the weighted error of Stump A.

\item
Calculate the weighted error of Stump B.

\item
Calculate the weighted error of Stump C.

\item
Which stump would AdaBoost select first?

\end{enumerate}
\begingroup\color{red}

\subsubsection*{Solution}
Initially,

\[
w_i=0.1
\]

for every patient.

Because all weights are equal, the weighted error is simply the fraction of incorrect predictions.

Therefore:

\[
{\epsilon_A=0.30}
\]

\[
{\epsilon_B=0.20}
\]

\[
{\epsilon_C=0.30}
\]

AdaBoost selects:

\[
{\text{Stump B: Blocked arteries}}
\]

because it has the lowest weighted error.

\par\endgroup

\clearpage
\section*{Part III - AdaBoost gives the stump a weight}
Once AdaBoost has selected a weak learner, it gives that learner an importance weight.

The weight is

\[
{
\alpha
=
\frac{1}{2}
\ln
\left(
\frac{1-\epsilon}{\epsilon}
\right)
}
\]

where $\epsilon$ is the weighted error of the stump.

Suppose the first stump has weighted error

\[
\epsilon=0.2.
\]

\subsubsection*{Exercise 3.1}
\begin{enumerate}
\item
Calculate $\alpha$.

\item
Is $\alpha$ positive or negative?

\item
What does a larger value of $\alpha$ mean?

\item
What would happen to $\alpha$ if the stump had an error of exactly 50\%?

\end{enumerate}
\begingroup\color{red}

\subsubsection*{Solution}
For the selected stump,

\[
\epsilon=0.20.
\]

The stump weight is

\[
\alpha=
\frac12
\ln\left(\frac{1-\epsilon}{\epsilon}\right).
\]

Therefore,

\[
\alpha=
\frac12
\ln\left(\frac{0.8}{0.2}\right)
\]

\[
=
\frac12\ln(4)
\]

\[
{\alpha\approx0.693}
\]

\subsubsection*{Interpretation}
The value is positive because the stump performs better than random guessing.

A \textbf{\textit{larger $\alpha$}} means that the weak learner has a smaller error and therefore receives more influence in the final ensemble.

If

\[
\epsilon=0.5,
\]

then

\[
\alpha=
\frac12\ln(1)=0.
\]

So a learner that performs no better than random guessing receives \textbf{\textit{zero weight}}.

\par\endgroup

\clearpage
\section*{Part IV - Updating the patient weights}
This is the key idea behind AdaBoost.

After the first stump has been trained:

\begin{itemize}
\item
Patients that were classified \textbf{\textit{correctly}} receive less weight.

\item
Patients that were classified \textbf{\textit{incorrectly}} receive more weight.

\end{itemize}
The update can be written as

\[
w_i'
=
w_i
\exp(-\alpha y_i h(x_i))
\]

where:

\begin{itemize}
\item
$y_i$ is the true label,

\item
$h(x_i)$ is the stump's prediction,

\item
$\alpha$ is the stump's weight.

\end{itemize}
If the prediction is correct:

\[
y_i h(x_i)=+1
\]

and therefore the weight is multiplied by

\[
e^{-\alpha}.
\]

If the prediction is incorrect:

\[
y_i h(x_i)=-1
\]

and therefore the weight is multiplied by

\[
e^{+\alpha}.
\]

The weights are then \textbf{\textit{normalized}} so that they add up to 1.

\subsection*{Task 4.1}
Suppose the first stump has

\[
\epsilon=0.2
\]

and therefore

\[
\alpha\approx0.693.
\]

Initially every patient has weight 0.1.

Calculate the new, unnormalized weight for:

\begin{itemize}
\item
a correctly classified patient;

\item
an incorrectly classified patient.

\end{itemize}
Use

\[
w_i'=w_i e^{-\alpha y_i h(x_i)}.
\]

Then normalize all weights.

\subsubsection*{Exercise 4.1}
\begin{enumerate}
\item
Which patients receive the largest weights?

\item
Why does AdaBoost deliberately give these patients larger weights?

\item
After this update, is every patient still equally important to the next stump?

\end{enumerate}
\begingroup\color{red}

\subsubsection*{Solution}
The first stump is:

Predict heart disease if blocked arteries = Yes.

It makes errors for:

\begin{itemize}
\item
Patient 2

\item
Patient 4

\end{itemize}
The other eight patients are correctly classified.

We have

\[
\alpha=0.693.
\]

For a correctly classified patient:

\[
w_i'=0.1e^{-0.693}
\approx0.1(0.5)
={0.05}.
\]

For an incorrectly classified patient:

\[
w_i'=0.1e^{0.693}
\approx0.1(2)
={0.20}.
\]

Thus the unnormalized weights are:

\begingroup
\small
\setlength{\tabcolsep}{4pt}
\begin{longtable}{l|l|l}
Patient & Correct? & Unnormalized weight \\[1.5pt]
\hline
\endfirsthead
Patient & Correct? & Unnormalized weight \\[1.5pt]
\hline
\endhead
1 & x & 0.05 \\[1.5pt]
2 & o & 0.20 \\[1.5pt]
3 & x & 0.05 \\[1.5pt]
4 & o & 0.20 \\[1.5pt]
5 & x & 0.05 \\[1.5pt]
6 & x & 0.05 \\[1.5pt]
7 & x & 0.05 \\[1.5pt]
8 & x & 0.05 \\[1.5pt]
9 & x & 0.05 \\[1.5pt]
10 & x & 0.05 \\[1.5pt]
\end{longtable}
\endgroup
The sum is

\[
8(0.05)+2(0.20)
=0.40+0.40
=0.80.
\]

We normalize by dividing every weight by 0.80.

Therefore:

\begin{itemize}
\item
Correctly classified patients:

\end{itemize}
\[
\frac{0.05}{0.80}
={0.0625}
\]

\begin{itemize}
\item
Incorrectly classified patients:

\end{itemize}
\[
\frac{0.20}{0.80}
={0.25}
\]

\subsubsection*{Final weights}
\begingroup
\small
\setlength{\tabcolsep}{4pt}
\begin{longtable}{l|l}
Patient & New weight \\[1.5pt]
\hline
\endfirsthead
Patient & New weight \\[1.5pt]
\hline
\endhead
1 & 0.0625 \\[1.5pt]
2 & \textbf{\textit{0.2500}} \\[1.5pt]
3 & 0.0625 \\[1.5pt]
4 & \textbf{\textit{0.2500}} \\[1.5pt]
5 & 0.0625 \\[1.5pt]
6 & 0.0625 \\[1.5pt]
7 & 0.0625 \\[1.5pt]
8 & 0.0625 \\[1.5pt]
9 & 0.0625 \\[1.5pt]
10 & 0.0625 \\[1.5pt]
\end{longtable}
\endgroup
The weights sum to 1.

\subsubsection*{Interpretation}
Patients \textbf{\textit{2 and 4}} now have much larger weights because the first stump failed to identify them.

The next stump therefore has a strong incentive to classify these patients correctly.

This is the central idea behind AdaBoost:

\textbf{\textit{Make difficult examples more important to subsequent learners.}}

\par\endgroup

\clearpage
\section*{Part V - The second stump}
The second stump is trained using the \textbf{\textit{new patient weights}}.

This means that making a mistake on a heavily weighted patient is more costly than making a mistake on a patient with a small weight.

Consider two possible second stumps:

\subsubsection*{Stump D}
\[
\text{If Chest pain = Yes}:
\quad \hat y=+1
\]

otherwise:

\[
\hat y=-1
\]

\subsubsection*{Stump E}
\[
\text{If Weight}>77.5\text{ kg}:
\quad \hat y=+1
\]

otherwise:

\[
\hat y=-1
\]

\subsubsection*{Exercise 5.1}
Using the updated patient weights from Part 4:

\begin{enumerate}
\item
Determine which patients Stump D classifies incorrectly.

\item
Calculate its \textbf{\textit{weighted error}}.

\item
Determine which patients Stump E classifies incorrectly.

\item
Calculate its \textbf{\textit{weighted error}}.

\item
Which stump should AdaBoost choose?

\item
Why might the stump with the smallest \emph{number} of mistakes not necessarily be the stump with the smallest \textbf{\textit{weighted error}}?

\end{enumerate}
\begingroup\color{red}

\subsubsection*{Solution}
\subsection*{Stump D - Chest pain}
Rule:

\[
\text{Chest pain = Yes}\Rightarrow+1
\]

It makes errors for patients:

\begin{itemize}
\item
Patient 3

\item
Patient 6

\item
Patient 9

\end{itemize}
Each has weight 0.0625.

Therefore:

\[
\epsilon_D
=
0.0625+0.0625+0.0625
\]

\[
{\epsilon_D=0.1875}
\]

\subsection*{Stump E - Weight}
Rule:

\[
\text{Weight}>77.5\text{ kg}\Rightarrow+1
\]

Predictions:

\begin{itemize}
\item
Patients 3, 4, 5, 6, 9, 10 $\to$ positive

\item
Patients 1, 2, 7, 8 $\to$ negative

\end{itemize}
Errors occur for:

\begin{itemize}
\item
Patient 2: infected, but predicted negative

\item
Patient 8: not infected, but predicted positive

\item
Patient 10: not infected, but predicted positive

\end{itemize}
Their weights are:

\[
w_2=0.25,\qquad
w_8=0.0625,\qquad
w_{10}=0.0625.
\]

Thus:

\[
\epsilon_E
=
0.25+0.0625+0.0625
\]

\[
{\epsilon_E=0.375}
\]

\subsubsection*{Which stump is selected?}
\[
\epsilon_D=0.1875
\]

versus

\[
\epsilon_E=0.375.
\]

Therefore AdaBoost selects:

\[
{\text{Stump D}}
\]

\subsubsection*{Important observation}
Stump D makes \textbf{\textit{three mistakes}}, while Stump E also makes three mistakes.

However, AdaBoost does \textbf{\textit{not}} treat all patients equally anymore.

Stump E makes a mistake on patient 2, who has a very large weight of 0.25.

Therefore its weighted error is much larger.

This demonstrates why AdaBoost uses \textbf{\textit{weighted error rather than simply counting mistakes}}.

\par\endgroup

\clearpage
\section*{Part VI - Combining the stumps}
AdaBoost does not simply take a majority vote where every stump has equal importance.

Instead, it uses the weighted vote

\[
{
F(x)
=
\alpha_1h_1(x)
+
\alpha_2h_2(x)
+
\alpha_3h_3(x)+\cdots
}
\]

The final prediction is

\[
{
\hat y=\operatorname{sign}(F(x))
}
\]

where:

\begin{itemize}
\item
$h_1,h_2,h_3,\ldots$ are the individual stumps;

\item
$\alpha_1,\alpha_2,\alpha_3,\ldots$ are their weights.

\end{itemize}
\subsection*{Example}
Suppose three stumps produce the following predictions for a new patient:

\begingroup
\small
\setlength{\tabcolsep}{4pt}
\begin{longtable}{l|l|l}
Stump & Prediction & Stump weight \\[1.5pt]
\hline
\endfirsthead
Stump & Prediction & Stump weight \\[1.5pt]
\hline
\endhead
$h_1$ & +1 & 0.8 \\[1.5pt]
$h_2$ & -1 & 0.4 \\[1.5pt]
$h_3$ & +1 & 0.3 \\[1.5pt]
\end{longtable}
\endgroup
Calculate

\[
F(x)
=
(0.8)(+1)+(0.4)(-1)+(0.3)(+1).
\]

\subsubsection*{Exercise 6.1}
\begin{enumerate}
\item
Calculate $F(x)$.

\item
What is the final AdaBoost prediction?

\item
Why does $h_1$ have more influence than $h_2$?

\end{enumerate}
\begingroup\color{red}

\subsubsection*{Solution}
Suppose three stumps make the following predictions:

\begingroup
\small
\setlength{\tabcolsep}{4pt}
\begin{longtable}{l|l|l}
Stump & Prediction & Weight \\[1.5pt]
\hline
\endfirsthead
Stump & Prediction & Weight \\[1.5pt]
\hline
\endhead
$h_1$ & +1 & 0.8 \\[1.5pt]
$h_2$ & -1 & 0.4 \\[1.5pt]
$h_3$ & +1 & 0.3 \\[1.5pt]
\end{longtable}
\endgroup
The combined score is

\[
F(x)
=
(0.8)(+1)
+
(0.4)(-1)
+
(0.3)(+1).
\]

Therefore,

\[
F(x)=0.8-0.4+0.3
\]

\[
{F(x)=0.7}
\]

Since the score is positive,

\[
{\hat y=+1}
\]

so the final model predicts \textbf{\textit{heart disease}}.

Stump $h_1$ has the largest influence because it has the largest $\alpha$.

\par\endgroup

\clearpage
\section*{Part VII - Understanding AdaBoost}
Complete the following sequence:

\[
\begin{gathered}
\text{Equal weights}\rightarrow\text{train stump}\rightarrow\text{calculate weighted error}\\
\rightarrow\text{calculate }\alpha\rightarrow\text{increase weights of mistakes}\rightarrow\text{train next stump}
\end{gathered}
\]

\subsubsection*{Exercise 7.1}
\begin{enumerate}
\item
Why does AdaBoost focus increasingly on difficult patients?

\item
What happens if a patient is repeatedly misclassified?

\item
Why are decision stumps useful as weak learners?

\item
Why can many simple stumps produce a much more powerful model when combined?

\end{enumerate}
\begingroup\color{red}

\subsubsection*{Solution}
\subsubsection*{1. Why does AdaBoost focus on difficult patients?}
Because incorrectly classified patients receive \textbf{\textit{larger weights}}.

This makes their classification errors more important when training the next weak learner.

\subsubsection*{2. What happens if a patient is repeatedly misclassified?}
Its weight can become increasingly important to subsequent learners.

The ensemble therefore increasingly focuses on difficult examples.

\subsubsection*{3. Why use decision stumps?}
A stump is a very simple and weak classifier.

For example:

``Does the patient have blocked arteries?''

It may not be accurate enough by itself, but many weak learners can be combined into a much stronger classifier.

\subsubsection*{4. Why can many simple stumps be powerful?}
Each stump can capture a different aspect of the data.

AdaBoost combines them using their individual weights, allowing the ensemble to make a more sophisticated overall decision.

\par\endgroup

\clearpage
\section*{Part VIII - From AdaBoost to Gradient Boosting}
AdaBoost and gradient boosting are both \textbf{\textit{boosting methods}}, but their mechanisms are different.

The central idea of gradient boosting is:

Build a model, examine its errors, and train the next model to improve those errors.

Instead of explicitly changing patient weights as AdaBoost does, gradient boosting can work with the \textbf{\textit{residuals}} or \textbf{\textit{negative gradients of a loss function}}.

\subsection*{A simplified regression example}
To understand the mechanism, temporarily imagine that we are predicting a continuous \textbf{\textit{heart-disease risk score}} rather than a yes/no diagnosis.

Suppose the initial model predicts the same risk for every patient:

\[
F_0(x)=0.5.
\]

For five patients, suppose the observed target values are:

\begingroup
\small
\setlength{\tabcolsep}{4pt}
\begin{longtable}{l|l|l}
Patient & Actual target $y$ & Initial prediction $F_0(x)$ \\[1.5pt]
\hline
\endfirsthead
Patient & Actual target $y$ & Initial prediction $F_0(x)$ \\[1.5pt]
\hline
\endhead
1 & 0.0 & 0.5 \\[1.5pt]
2 & 1.0 & 0.5 \\[1.5pt]
3 & 1.0 & 0.5 \\[1.5pt]
4 & 0.0 & 0.5 \\[1.5pt]
5 & 1.0 & 0.5 \\[1.5pt]
\end{longtable}
\endgroup
For squared-error loss, the residual is

\[
r_i=y_i-F_0(x_i).
\]

\subsubsection*{Exercise 8.1}
\begin{enumerate}
\item
Calculate the residual for each patient.

\item
Which patients does the current model underestimate?

\item
Which patients does it overestimate?

\item
What should the next decision tree try to predict?

\end{enumerate}
\begingroup\color{red}

\subsubsection*{Solution}
The simplified example uses

\[
F_0(x)=0.5.
\]

The residual is

\[
r_i=y_i-F_0(x_i).
\]

\subsubsection*{Residuals}
\begingroup
\small
\setlength{\tabcolsep}{4pt}
\begin{longtable}{l|l|l|l}
Patient & Actual $y$ & Prediction $F_0$ & Residual $r=y-F_0$ \\[1.5pt]
\hline
\endfirsthead
Patient & Actual $y$ & Prediction $F_0$ & Residual $r=y-F_0$ \\[1.5pt]
\hline
\endhead
1 & 0.0 & 0.5 & -0.5 \\[1.5pt]
2 & 1.0 & 0.5 & +0.5 \\[1.5pt]
3 & 1.0 & 0.5 & +0.5 \\[1.5pt]
4 & 0.0 & 0.5 & -0.5 \\[1.5pt]
5 & 1.0 & 0.5 & +0.5 \\[1.5pt]
\end{longtable}
\endgroup
Therefore:

\[
{
r=(-0.5,+0.5,+0.5,-0.5,+0.5)
}
\]

Patients 2, 3, and 5 have positive residuals, meaning the model is \textbf{\textit{underestimating}} their target.

Patients 1 and 4 have negative residuals, meaning the model is \textbf{\textit{overestimating}} their target.

The next tree therefore tries to learn these residuals - in other words, it tries to learn how the current model should be \textbf{\textit{corrected}}.

\par\endgroup

\clearpage
\section*{Part IX - Gradient boosting with a decision stump}
Suppose the next stump learns the following rule:

\[
\text{If Age}>50:
\quad r=+0.3
\]

otherwise:

\[
r=-0.2.
\]

The new model is obtained by \textbf{\textit{adding the correction}} to the previous model:

\[
{
F_1(x)=F_0(x)+\eta h_1(x)
}
\]

where $\eta$ is the \textbf{\textit{learning rate}}.

For this exercise, assume

\[
\eta=1.
\]

\subsubsection*{Exercise 9.1}
\begin{enumerate}
\item
Calculate the new prediction for a patient aged 60.

\item
Calculate the new prediction for a patient aged 40.

\item
Explain why the second tree is called a \textbf{\textit{correction}} to the first model.

\item
What would happen if we continued adding more trees?

\end{enumerate}
\begingroup\color{red}

\subsubsection*{Solution}
The second stump predicts:

\[
\text{Age}>50
\Rightarrow r=+0.3
\]

otherwise:

\[
r=-0.2.
\]

The new model is

\[
F_1(x)=F_0(x)+\eta h_1(x).
\]

Given

\[
F_0=0.5
\]

and

\[
\eta=1,
\]

\subsubsection*{Patient aged 60}
Since

\[
60>50,
\]

the stump predicts

\[
+0.3.
\]

Therefore:

\[
F_1=0.5+0.3
\]

\[
{F_1=0.8}
\]

\subsubsection*{Patient aged 40}
Since

\[
40<50,
\]

the stump predicts

\[
-0.2.
\]

Therefore:

\[
F_1=0.5-0.2
\]

\[
{F_1=0.3}
\]

The second tree is a \textbf{\textit{correction}} because it does not replace the first model. It adds information that attempts to reduce the existing prediction errors.

With additional trees:

\[
F_2(x)=F_1(x)+\eta h_2(x)
\]

\[
F_3(x)=F_2(x)+\eta h_3(x)
\]

and so on.

The model gradually becomes more sophisticated.

\par\endgroup

\clearpage
\section*{Part X - AdaBoost vs. Gradient Boosting}
\subsubsection*{Exercise 10.1}
Complete the comparison table.

\begingroup
\small
\setlength{\tabcolsep}{4pt}
\begin{longtable}{l|l|l}
\rule{0pt}{2.4ex} & AdaBoost & Gradient Boosting \\[3pt]
\hline
\endfirsthead
\rule{0pt}{2.4ex} & AdaBoost & Gradient Boosting \\[3pt]
\hline
\endhead
\rule{0pt}{2.4ex}Basic building block &  $\phantom{-0.00}$ &  $\phantom{-0.00}$ \\[3pt]
\rule{0pt}{2.4ex}What happens after an error? &  $\phantom{-0.00}$ &  $\phantom{-0.00}$ \\[3pt]
\rule{0pt}{2.4ex}How is the next learner guided? &  $\phantom{-0.00}$ &  $\phantom{-0.00}$ \\[3pt]
\rule{0pt}{2.4ex}How are learners combined? &  $\phantom{-0.00}$ &  $\phantom{-0.00}$ \\[3pt]
\rule{0pt}{2.4ex}Main idea &  $\phantom{-0.00}$ &  $\phantom{-0.00}$ \\[3pt]
\end{longtable}
\endgroup
\begingroup\color{red}

\subsubsection*{Solution}
\begingroup
\small
\setlength{\tabcolsep}{4pt}
\begin{longtable}{l|l|l}
\rule{0pt}{2.4ex} & AdaBoost & Gradient Boosting \\[3pt]
\hline
\endfirsthead
\rule{0pt}{2.4ex} & AdaBoost & Gradient Boosting \\[3pt]
\hline
\endhead
Basic building block & Weak learners such as stumps & Usually shallow decision trees \\[1.5pt]
What happens after an error? & \shortstack[l]{Misclassified examples receive \\ larger weights} & \shortstack[l]{The next learner focuses on the \\ current model's residuals/negative \\ gradients} \\[1.5pt]
How is the next learner guided? & By the updated example weights & \shortstack[l]{By the gradient of the loss \\ function} \\[1.5pt]
How are learners combined? & Weighted sum/vote using $\alpha$ & \shortstack[l]{Sequentially added to the existing \\ model} \\[1.5pt]
Main idea & \shortstack[l]{Focus more strongly on difficult \\ training examples} & \shortstack[l]{Add models that correct the \\ current model's errors} \\[1.5pt]
\end{longtable}
\endgroup
\par\endgroup

\clearpage
\section*{Part XI - Short conceptual questions}
\subsubsection*{Exercise 11.1}
\begin{enumerate}
\item
In your own words, explain the difference between \textbf{\textit{a stump}} and an \textbf{\textit{ensemble of stumps}}.

\item
Explain why AdaBoost changes the weights of training examples.

\item
Explain how the stump's $\alpha$ determines its influence.

\item
Explain how gradient boosting uses the errors of the current model.

\item
What is the main conceptual similarity between AdaBoost and gradient boosting?

\item
What is the main conceptual difference?

\end{enumerate}
\begingroup\color{red}

\subsubsection*{Solution}
\subsubsection*{1. What is the difference between a stump and an ensemble of stumps?}
A \textbf{\textit{stump}} is a single, very simple decision rule based on one split.

An \textbf{\textit{ensemble of stumps}} combines many such rules to produce a more powerful classifier.

\subsubsection*{2. Why does AdaBoost change training-example weights?}
To make incorrectly classified examples more important to subsequent learners.

The next stump therefore concentrates more on the cases that previous stumps struggled with.

\subsubsection*{3. How does $\alpha$ determine a stump's influence?}
\[
\alpha=
\frac12\ln\left(\frac{1-\epsilon}{\epsilon}\right)
\]

A stump with a smaller error $\epsilon$ receives a larger $\alpha$, giving it greater influence in the final prediction.

\subsubsection*{4. How does gradient boosting use errors?}
For a given loss function, gradient boosting determines the direction in which the current predictions should be changed to reduce the loss.

A new tree is trained to approximate this \textbf{\textit{correction}}, and the correction is added to the existing model.

\subsubsection*{5. What is the main similarity between AdaBoost and gradient boosting?}
Both methods build a strong model \textbf{\textit{sequentially from multiple weak learners}}.

Each new learner attempts to improve the ensemble produced by the previous learners.

\subsubsection*{6. What is the main difference?}
The key difference is \textbf{\textit{how the next learner is determined}}:

\begin{itemize}
\item
\textbf{\textit{AdaBoost:}} changes the weights of training examples, emphasizing examples that were previously misclassified.

\item
\textbf{\textit{Gradient boosting:}} fits the next learner to the \textbf{\textit{negative gradient of the loss}} (often introduced as residuals for squared-error regression).

\end{itemize}
\par\endgroup

\section*{Summary}
The important ideas to take away are:

\subsubsection*{Decision stump}
A very simple decision tree containing only \textbf{\textit{one decision}}.

\subsubsection*{AdaBoost}
AdaBoost repeatedly:

\begin{enumerate}
\item
Starts with equal patient weights.

\item
Trains a weak learner.

\item
Measures its weighted error.

\item
Gives the learner a weight $\alpha$.

\item
Increases the weights of incorrectly classified patients.

\item
Trains another learner that focuses more on difficult cases.

\item
Combines the learners using their weights.

\end{enumerate}
\subsubsection*{Gradient boosting}
Gradient boosting repeatedly:

\begin{enumerate}
\item
Starts with a simple prediction.

\item
Calculates how the current model is wrong according to a loss function.

\item
Trains a new tree to predict a \textbf{\textit{correction}} to the current model.

\item
Adds that correction to the existing model.

\item
Repeats the process.

\end{enumerate}
Thus, both methods turn many \textbf{\textit{weak learners}} into a stronger model, but they determine what the next learner should focus on in different ways.


\end{document}
